\documentclass[10pt,a4paper]{amsart}
\usepackage[margin=19mm]{geometry}
\usepackage{amsmath,amssymb,mathtools}
\usepackage[scr=rsfs,cal=euler]{mathalfa}
\usepackage{xfrac}
\usepackage[unicode,hidelinks,bookmarks=false]{hyperref}
\newcommand{\ie}{i.e.\ }
\newcommand{\cf}{cf.\ }
\newcommand{\resp}{respectively\ }
\newcommand{\Subsection}{\subsection}
\providecommand{\nobreakdash}{\nobreak\hbox{-}}
\setlength{\emergencystretch}{2em}
\multlinegap0pt
\usepackage{amssymb}
\newtheorem{theorem}{Theorem}
\title{Generic three-parameter unfolding of a completely degenerate germ of a vector field on $\mathbb{R}^3$}
\author{Samuel Diangitukulu Ndimba, Aimable Sayinzoga, Bertuel Tangue Ndawa, Jonathan Ibia Kobongo, Joseph Mpulu Kikwasini}
\date{Source: arXiv:2608.02074v2 (CC BY 4.0)}
\begin{document}
\maketitle

\noindent\emph{Source and licence.} Generic three-parameter unfolding of a completely degenerate germ of a vector field on $\mathbb{R}^3$. Version arXiv:2608.02074v2 (CC BY 4.0), distributed by arXiv under the Creative Commons Attribution 4.0 International licence, http://creativecommons.org/licenses/by/4.0/, as recorded in the arXiv licence metadata for this exact version and shown on the abstract page https://arxiv.org/abs/2608.02074v2. Full licence notice: \texttt{licenses/arxiv-2608.02074-CC-BY-4.0.txt}.\par\smallskip

\noindent\textit{Editorial scope.} Counted unit: the article's theorem theorem (source0.tex lines 1076--1157). The article's complete original proof of it is delivered with it (source0.tex lines 1160--1293, one \texttt{proof} environment of 134 lines). No mathematics is written, repaired or re-derived: the statement and the proof are the article's own lines, in the article's own notation, and no retained line is substituted or re-flowed.  Retained uncounted as the source-local context the proof needs: lines 337--342; lines 349--353; lines 782--808.  Nothing was omitted from any retained span, so the package carries no omission record.  No retained line contains a citation, so the wrapper carries no bibliography and no bibliography entry.  No retained line contains a figure environment, an image-inclusion command or drawing code, so no image is delivered and the package declares no asset dependency.  Editorial formatting only, in addition: an AMS \texttt{amsart} preamble that restates the article's own declarations and macros used by the retained lines, namely the environments \texttt{theorem} on separate counters, together with the standard packages those lines need.  The retained environments print in this document as Theorem 1 in order of appearance, whereas the article prints the same items with the numbering of its own full document.  The complete native source of the article as distributed in the arXiv e-print for this exact version is bundled unchanged as \texttt{sources/206666/source0.tex}.
\begin{equation}\label{GSUH1}\tag{H1}
\begin{array}{l}
X_0=x_2\frac{\partial}{\partial x_1}+x_3\frac{\partial}{\partial x_2} + f_0 (x)\frac{\partial}{\partial x_3},
 \\f_0(x) = \sum_{i,j=1}^{3} x_i x_j a_{ij}(x),\; (a_{ij})\in C^m(M, S_n(\mathbb{R})),\mbox{ and } a_{11}(0)=1;
\end{array}
\end{equation}

\begin{equation}\label{GSUH3}\tag{H3}
\begin{array}{l}f (\lambda , x)+ f_3 (\lambda , x) = f_{30} (\lambda)+x_1 f_{31}(\lambda ) + x_{1}^{2}f_{32}(\lambda ,x_1)+x_2 h_1 (\lambda , x)\\
\hspace{3.6cm}+ x_3 h_2 (\lambda ,x ),\\ 
f_{3i},\; h_1\; h_2\in C^m\left(\mathbb{R}^k\times\mathbb{R}^3\right),\;  i=1,2,3,\;f_{32} (0,0)=1,\\  f_{30}(0,0) = f_{3 1}(0,0) =h_i (0,0),\; i=1,2\end{array}
\end{equation}

\begin{equation}
	\label{GSUH5}
	\tag{H5}
	\begin{array}{l}
		X(\lambda,x)=X_{\lambda}(x)
		=
		x_2\frac{\partial}{\partial x_1}
		+
		x_3\frac{\partial}{\partial x_2}
		+
		f_{\lambda}(x)\frac{\partial}{\partial x_3},
		\\[0.2cm]
		\begin{cases}
			f_{\lambda}(x)
			=
			\lambda_1+\lambda_2x_2+\lambda_3x_3+x_1^2
			+
			\displaystyle\sum_{i,j=1}^{3}x_ix_jA_{ij}(\lambda,x),
			\\[0.2cm]
			A_{ij}=A_{ji},
			\qquad
			A_{ij}\in C^m(\mathbb{R}^3\times \mathbb{R}^3),
			\\[0.2cm]
			A_{11}(0,0)=A_{22}(0,0)=A_{23}(0,0)=A_{33}(0,0)=0.
		\end{cases}
	\end{array}
\end{equation}

\begin{theorem}
	Let $X\in \mathfrak{X}^{m}(\mathbb{R}^{3}\mid \mathbb{R}^{3})$
	satisfy hypotheses \eqref{GSUH1}--\eqref{GSUH3}. Let
	$J(\lambda,x)$ denote the Jacobian matrix of $X_\lambda$ with respect to
	$x$. If $r_1\in \mathbb{R}$, $r_2=\alpha+i\beta$, $r_3=\alpha-i\beta$
	are the eigenvalues of $J(\lambda,x)$, then the germ
	$(X_\lambda,0)$ possesses the following local bifurcations.
	
	\begin{enumerate}
		\item Three singularities of codimension one.
		
		\begin{enumerate}
			\item The singularity $	x^1=(0,0,0)$
			appears on the surface
			$$
			S_1=
			\left\{
			\lambda\in \mathbb{R}^3:
			\lambda_1=0,\;
			\lambda_3\neq 0,\;
			4\lambda_2+\lambda_3^2\leq 0
			\right\}.
			$$
			The eigenvalues are $r_1=0$,	and
			$$
			r_{2,3}
			=
			\frac{1}{2}
			\left(
			\lambda_3
			\pm
			i\left|4\lambda_2+\lambda_3^2\right|^{1/2}
			\right).
			$$
			
			\item For $\lambda_1<0$, let $v_\varepsilon	=-\varepsilon\sqrt{-\lambda_1}$, $\varepsilon=\pm 1,$
			be the corresponding solution, up to higher-order terms, of
			$$
			\lambda_1+x_1^2
			\left(
			1+A_{11}(\lambda,x_1)
			\right)=0.
			$$
			Then $x^\varepsilon=(v_\varepsilon,0,0)$, $\varepsilon=\pm 1,$
			are singularities associated with the surface
			$$
			S_2^\varepsilon
			=
			\left\{
			\lambda:
			\gamma_{2\varepsilon}<0,\;
			\gamma_{3\varepsilon}\neq 0,\;
			\gamma_{1\varepsilon}
			+
			\gamma_{2\varepsilon}\gamma_{3\varepsilon}=0
			\right\},
			$$
			where
			$$
			\gamma_{i\varepsilon}
			=
			\frac{\partial f_\lambda}{\partial x_i}
			(\lambda,x^\varepsilon),
			\qquad
			i=1,2,3.
			$$
			The eigenvalues are $r_1=\gamma_{3\varepsilon},$ $r_2=-r_3=i\sqrt{-\gamma_{2\varepsilon}}.$
		\end{enumerate}
		
		\item Two singularities of codimension two.
		
		\begin{enumerate}
			\item For $\lambda\in C_1=\{0\}\times \mathbb{R}_{-}^{\ast}\times \{0\},$
			the singularity $x^2=(0,0,0)$
			has eigenvalues $r_1=0$, $r_2=-r_3=i\sqrt{-\lambda_2}.$
			
			\item For $\lambda\in C_2=\{0\}\times \{0\}\times \mathbb{R}^{\ast},$
			the singularity $x^3=(0,0,0)$ has eigenvalues $r_1=\lambda_3,$ 
			$r_2=r_3=0.$
		\end{enumerate}
	\end{enumerate}
\end{theorem}

\begin{proof}
	The singularities of $X_\lambda$ are the solutions of
	$$
	x_2=0,
	\qquad
	x_3=0,
	\qquad
	f_\lambda(x_1,0,0)=0.
	$$
	Using \eqref{GSUH5}, this last equation becomes
	$$
	\lambda_1
	+
	x_1^2
	\left(
	1+A_{11}(\lambda,x_1,0,0)
	\right)
	=0.
	$$
	
	At a singularity $x^\ast=(v,0,0)$, the Jacobian matrix of $X_\lambda$ is
	$$
	J(\lambda,x^\ast)
	=
	\begin{pmatrix}
		0 & 1 & 0\\
		0 & 0 & 1\\
		\gamma_1 & \gamma_2 & \gamma_3
	\end{pmatrix},
	$$
	where
	$$
	\gamma_i
	=
	\frac{\partial f_\lambda}{\partial x_i}
	(\lambda,x^\ast),
	\qquad
	i=1,2,3.
	$$
	The characteristic polynomial is therefore
	$$
	P(r)
	=
	r^3-\gamma_3r^2-\gamma_2r-\gamma_1.
	$$
	
	At the origin $x^1=(0,0,0)$, when $\lambda_1=0$, one has
	$$
	\gamma_1=0,
	\qquad
	\gamma_2=\lambda_2,
	\qquad
	\gamma_3=\lambda_3.
	$$
	Thus
	$$
	P(r)
	=
	r(r^2-\lambda_3r-\lambda_2).
	$$
	Hence one eigenvalue is $r_1=0,$
	and the other two are
	$$
	r_{2,3}
	=
	\frac{1}{2}
	\left(
	\lambda_3
	\pm
	\sqrt{\lambda_3^2+4\lambda_2}
	\right).
	$$
	If $\lambda_3^2+4\lambda_2\leq 0,$
	then this pair is complex conjugate and may be written as
	$$
	r_{2,3}
	=
	\frac{1}{2}
	\left(
	\lambda_3
	\pm
	i\left|4\lambda_2+\lambda_3^2\right|^{1/2}
	\right).
	$$
	This gives the first codimension-one case.
	
	Now assume $\lambda_1<0$. Then the equation
	$$
	\lambda_1
	+
	x_1^2
	\left(
	1+A_{11}(\lambda,x_1,0,0)
	\right)
	=0
	$$
	has two nearby solutions
	$$
	v_\varepsilon
	=
	-\varepsilon\sqrt{-\lambda_1}
	+
	\text{higher-order terms},
	\qquad
	\varepsilon=\pm 1.
	$$
	Thus there are two singularities $x^\varepsilon=(v_\varepsilon,0,0).$
	At such a point, the eigenvalues are the roots of
	$$
	r^3-\gamma_{3\varepsilon}r^2
	-\gamma_{2\varepsilon}r
	-\gamma_{1\varepsilon}
	=0.
	$$
	If this polynomial has one real eigenvalue $r_1=\gamma_{3\varepsilon}$
	and a purely imaginary pair $r_{2,3}=\pm i\sqrt{-\gamma_{2\varepsilon}},$
	then necessarily $\gamma_{2\varepsilon}<0$
	and $\gamma_{1\varepsilon}+\gamma_{2\varepsilon}\gamma_{3\varepsilon}=0.$
	This gives the surfaces $S_2^\varepsilon$.
	
	Finally, the codimension-two cases occur when degeneracies in the previous
	codimension-one surfaces overlap.
	
	If $\lambda_1=0,$ $\lambda_3=0,$ $\lambda_2<0,$
	then at the origin the characteristic polynomial is $
	P(r)=r(r^2-\lambda_2),$
	and therefore $r_1=0,$ $r_{2,3}=\pm i\sqrt{-\lambda_2}.$
	This gives the curve $C_1$.
	
	If $\lambda_1=0,$ $\lambda_2=0,$ $\lambda_3\neq 0,$
	then $P(r)=r^2(r-\lambda_3),$
	so that $r_1=\lambda_3,$ $r_2=r_3=0.$
	This gives the curve $C_2$.
\end{proof}

\end{document}
